\documentclass[11pt,openany]{book}

% ============================================================
%  preamble.tex — shared preamble for the algorithm reference.
%  Typographic conventions inspired by the canonical algorithms
%  textbook: Computer Modern, two-column-friendly math, boxed
%  pseudocode with line numbers, margin-callout invariants.
% ============================================================

\usepackage[T1]{fontenc}
\usepackage[utf8]{inputenc}
\usepackage{lmodern}
\usepackage{textcomp}
\usepackage[margin=1.25in]{geometry}
\usepackage{amsmath,amssymb,amsthm}
\usepackage{mathtools}
\usepackage{graphicx}
\usepackage{xcolor}
\usepackage{microtype}
\usepackage{enumitem}
\usepackage{booktabs}
\usepackage{tabularx}
\usepackage{array}
\usepackage{caption}
\usepackage[nottoc,notlot,notlof]{tocbibind}
\usepackage{titlesec}
\usepackage{fancyhdr}
\usepackage{framed}
\usepackage[ruled,linesnumbered,noend]{algorithm2e}
\usepackage{hyperref}
\hypersetup{
    colorlinks=true,
    linkcolor=black,
    citecolor=black,
    urlcolor=black,
    pdftitle={GraphHunter Algorithms Reference},
    pdfauthor={GraphHunter Team}
}

% ------------------------------------------------------------
% Page style — running header with chapter + section, as in
% the reference algorithms text.
% ------------------------------------------------------------
\pagestyle{fancy}
\fancyhf{}
\renewcommand{\headrulewidth}{0.4pt}
\fancyhead[LE,RO]{\thepage}
\fancyhead[LO]{\small\itshape\nouppercase{\rightmark}}
\fancyhead[RE]{\small\itshape\nouppercase{\leftmark}}

% ------------------------------------------------------------
% Chapter / section styling — restrained, serif, no color.
% ------------------------------------------------------------
\titleformat{\chapter}[display]
  {\normalfont\Huge\bfseries}
  {\chaptertitlename\ \thechapter}{20pt}{\Huge}
\titleformat*{\section}{\Large\bfseries}
\titleformat*{\subsection}{\large\bfseries}
\titleformat*{\subsubsection}{\normalsize\bfseries\itshape}

% ------------------------------------------------------------
% Theorem-like environments with custom styles matching the
% "Theorem / Lemma / Invariant" layout.
% ------------------------------------------------------------
\newtheoremstyle{ghtheorem}
  {12pt}{12pt}{\itshape}{}{\bfseries}{.}{.5em}
  {\thmname{#1}\thmnumber{ #2}\thmnote{ (#3)}}
\newtheoremstyle{ghdefinition}
  {12pt}{12pt}{}{}{\bfseries}{.}{.5em}
  {\thmname{#1}\thmnumber{ #2}\thmnote{ (#3)}}

\theoremstyle{ghtheorem}
\newtheorem{theorem}{Theorem}[chapter]
\newtheorem{lemma}[theorem]{Lemma}
\newtheorem{corollary}[theorem]{Corollary}
\newtheorem{invariant}[theorem]{Invariant}

\theoremstyle{ghdefinition}
\newtheorem{definition}[theorem]{Definition}
\newtheorem{problem}[theorem]{Problem}
\newtheorem{remark}[theorem]{Remark}

% ------------------------------------------------------------
% Algorithm2e tuning: numbered lines, bold keywords, en-dash
% comments — matches the pseudocode convention of the canonical
% text. Keywords small-caps, variable names italic math.
% ------------------------------------------------------------
\SetKwSty{textsc}
\SetFuncSty{textsc}
\SetArgSty{textit}
\SetKwInOut{Input}{Input}
\SetKwInOut{Output}{Output}
\SetKwInOut{Invariant}{Invariant}
\SetKwFor{For}{for}{do}{end for}
\SetKwFor{While}{while}{do}{end while}
\SetKw{KwReturn}{return}
\SetKw{KwAnd}{and}
\SetKw{KwOr}{or}
\SetKw{KwNot}{not}
\SetKwComment{Comment}{\hfill$\triangleright$\ }{}
\DontPrintSemicolon
\LinesNumbered

% ------------------------------------------------------------
% Custom math operators and shorthands used across chapters.
% ------------------------------------------------------------
\DeclareMathOperator*{\argmax}{arg\,max}
\DeclareMathOperator*{\argmin}{arg\,min}
\DeclareMathOperator{\freq}{freq}
\DeclareMathOperator{\degout}{deg}
\newcommand{\bigO}{\mathcal{O}}
\newcommand{\N}{\mathbb{N}}
\newcommand{\Z}{\mathbb{Z}}
\newcommand{\R}{\mathbb{R}}
\newcommand{\eps}{\varepsilon}
\newcommand{\set}[1]{\{#1\}}
\newcommand{\abs}[1]{\lvert #1 \rvert}
\newcommand{\card}[1]{\lvert #1 \rvert}
\newcommand{\code}[1]{\texttt{#1}}
\newcommand{\unit}[1]{\ensuremath{\,\mathrm{#1}}}

% ------------------------------------------------------------
% "Implementation note" side-block: a lightly framed aside that
% connects the formal specification to the code that realizes it.
% ------------------------------------------------------------
\newenvironment{implnote}[1][Implementation note]{%
  \par\medskip\noindent
  \begin{framed}\small
  \textbf{#1.}\ \ignorespaces
}{\end{framed}\par\medskip}

% End of preamble.


\title{\bfseries GraphHunter\\[0.6em]\Large Algorithms Reference}
\author{GraphHunter Engineering}
\date{\today}

\begin{document}

\frontmatter
\maketitle

\begin{center}
\vspace*{2em}
\small
\textit{A specification-style reference for the three numerical and
graph routines at the core of GraphHunter: temporal depth-first
search, five-component path anomaly scoring, and SIMD-accelerated
sorted set intersection.}
\end{center}

\cleardoublepage
\tableofcontents

\mainmatter

% ============================================================
\chapter{Preliminaries}
\label{chap:preliminaries}

This document specifies three algorithms used in GraphHunter's core
detection pipeline. Each chapter follows the same structure: a
\textbf{problem statement} formalizing the input and output, an
\textbf{algorithm} description in pseudocode, a \textbf{complexity}
analysis, a list of \textbf{invariants} the implementation must
satisfy, and \textbf{references} to the literature from which the
techniques were drawn.

The documents are \emph{specifications}, not code walkthroughs. Line
numbers and function names drift as the codebase evolves; the math
and the invariants do not. Each invariant is pinned by at least one
property-based test in the Rust test suite; those tests are identified
by their file name inside \texttt{core/graph-engine/tests/} or
\texttt{core/matcher-ffi/tests/}.

\section*{Notation}

Throughout, $G = (V, E)$ denotes a directed, labeled, time-stamped
graph. Each node $v \in V$ carries an entity type $\tau(v)$ drawn
from a finite alphabet. Each edge $e = (u, v, r, t) \in E$ carries
a source $u$, destination $v$, relation-type label $r$, and an
integer timestamp $t \in \Z$.

A \emph{hypothesis} $H = (s_1, s_2, \dots, s_k)$ is an ordered
sequence of \emph{steps}; each step $s_i = (\tau^{\mathrm{o}}_i,
r_i, \tau^{\mathrm{d}}_i, P_i)$ specifies an origin type, a
relation type, a destination type, and a set of predicates on
edge and destination metadata.

A \emph{path} $\pi = v_0 \to v_1 \to \dots \to v_k$ is a sequence
of $k+1$ nodes connected by $k$ edges; path indexing is zero-based.

% ============================================================
\chapter{Temporal Depth-First Search}
\label{chap:dfs}

\section{Problem}
\label{sec:dfs:problem}

\begin{problem}[Hypothesis matching with temporal causality]
\label{prob:dfs}
Given a graph $G$, a validated hypothesis $H = (s_1, \dots, s_k)$,
an optional time window $[t_{\text{start}}, t_{\text{end}}]$, and a
result cap $c \in \N$, enumerate up to $c$ paths
$\pi = v_0 \to v_1 \to \dots \to v_k$ satisfying:
\begin{enumerate}
  \item \textbf{Type chain.} For every $i \in [1, k]$,
        $\tau(v_{i-1}) = \tau^{\mathrm{o}}_i$ and
        $\tau(v_i) = \tau^{\mathrm{d}}_i$. A wildcard type
        (\textsc{Any}) matches all.
  \item \textbf{Relation-type match.} There exists an edge
        $(v_{i-1}, v_i, r, t_i) \in E$ with $r = r_i$, or
        $r_i = \textsc{Any}$.
  \item \textbf{Temporal causality.} $t_1 \le t_2 \le \dots \le t_k$.
  \item \textbf{Time-window bound.} If the window is set, every
        $t_i \in [t_{\text{start}}, t_{\text{end}}]$.
  \item \textbf{$k$-simplicity.} The nodes $v_0, v_1, \dots, v_k$
        are pairwise distinct.
  \item \textbf{Predicate match.} The edge and destination metadata
        satisfy every predicate in $P_i$.
\end{enumerate}
\end{problem}

\section{Algorithm}
\label{sec:dfs:algorithm}

The matcher performs depth-first search with backtracking,
parameterized by the current hypothesis step. The pseudocode below
extends the standard DFS (Cormen \emph{et al.}~\cite{clrs} \S22.3)
with a pattern pointer, a temporal guard, and predicate evaluation.

\begin{algorithm}[H]
\caption{\textsc{Search-Temporal-Pattern}$(G, H, [t_{\text{lo}}, t_{\text{hi}}], c)$}
\label{alg:dfs}
\Input{graph $G$; hypothesis $H$; window $[t_{\text{lo}}, t_{\text{hi}}]$; cap $c$}
\Output{set $R$ of matching paths, with $\abs{R} \le c$; truncation flag}

$R \gets \emptyset$; $\mathit{trunc} \gets \bot$\;
$S \gets \set{v \in V : \tau(v) = \tau^{\mathrm{o}}_1 \lor \tau^{\mathrm{o}}_1 = \textsc{Any}}$\;
\For{each $v \in S$}{
  \If{$\abs{R} \ge c$}{$\mathit{trunc} \gets \top$; \textbf{break}\;}
  \textsc{Dfs}$(v, 1, [v], \set{v})$\;
}
\KwReturn $(R, \mathit{trunc})$\;

\BlankLine
\SetKwFunction{Dfs}{Dfs}
\SetKwProg{myproc}{Procedure}{}{}
\myproc{\Dfs{$u, i, \pi, V_{\mathrm{seen}}$}}{
  \If{$i > k$}{
    $R \gets R \cup \set{\pi}$\;
    \KwReturn\;
  }
  \For{each edge $(u, w, r, t) \in E$ in ascending $t$ order}{
    \If{$\tau(w) \ne \tau^{\mathrm{d}}_i \land \tau^{\mathrm{d}}_i \ne \textsc{Any}$}{\textbf{continue}\;}
    \If{$r \ne r_i \land r_i \ne \textsc{Any}$}{\textbf{continue}\;}
    \If{$w \in V_{\mathrm{seen}}$}{\textbf{continue} \Comment*[r]{$k$-simplicity}}
    \If{$\pi$ non-empty \textbf{and} $t < t_{\mathrm{last}}$}{\textbf{continue}\;}
    \If{window set \textbf{and} $t \notin [t_{\text{lo}}, t_{\text{hi}}]$}{\textbf{continue}\;}
    \If{predicates $P_i$ reject $(e, w)$}{\textbf{continue}\;}
    \textsc{Dfs}$(w, i+1, \pi \cdot w, V_{\mathrm{seen}} \cup \set{w})$\;
  }
}
\end{algorithm}

The driver selects candidate start nodes by type, then launches
\textsc{Dfs} from each. The procedure emits on success ($i > k$),
otherwise recurses on every edge that satisfies the combined type,
temporal, and predicate constraints. Backtracking unwinds the
\emph{visited} set and the partial path in LIFO order.

\section{Complexity}
\label{sec:dfs:complexity}

Let $n = \abs{V}$, $m = \abs{E}$, $k = \abs{H}$, and
$b = \max_{v \in V} \degout_{\text{out}}(v)$.

\begin{theorem}[Worst-case running time]
\label{thm:dfs:worst}
\textsc{Search-Temporal-Pattern} runs in $\bigO(n \cdot b^k)$ time
and uses $\bigO(k)$ stack depth plus $\bigO(n)$ memory for the
visited set.
\end{theorem}

\begin{proof}[Sketch]
The outer loop over candidate start nodes is $\bigO(n)$. Each
\textsc{Dfs} call at depth $i$ iterates at most $b$ edges and
recurses at most $b$ times at depth $i+1$. Unfolding the recursion
to depth $k$ yields $b^k$ work per start node. The stack depth is
bounded by $k$, and \emph{visited} has $\bigO(n)$ entries.
\end{proof}

In practice, the type filter and the causality constraint reduce the
effective branching factor to $b' \ll b$. On the \texttt{User-Auth-IP}
benchmark with $n = 10{,}000$, observed wall time is $\approx 15\unit{ms}$,
consistent with $b' \cdot k \approx 30$ edge visits per start node.

\begin{remark}[Early termination]
The enumeration stops once $\abs{R} = c$, at which point the caller
receives $\mathit{trunc} = \top$. This caps wall time at the cost of
result completeness; downstream code must treat the result set as a
sample, not a census.
\end{remark}

\section{Invariants}
\label{sec:dfs:invariants}

The following invariants must hold for every call; violations are
bugs, not edge cases.

\begin{invariant}[Hypothesis well-formedness]
\label{inv:dfs:hypo}
$H$ has passed \textsc{Hypothesis.validate} before
\textsc{Search-Temporal-Pattern} is invoked. In particular,
$k \ge 1$ and $\tau^{\mathrm{d}}_i = \tau^{\mathrm{o}}_{i+1}$ for
every $i \in [1, k-1]$.
\end{invariant}

\begin{invariant}[Visited-set well-balance]
\label{inv:dfs:visited}
Every insertion into $V_{\mathrm{seen}}$ is matched by a deletion
along the backtracking path. In particular, on return to the driver
loop, $V_{\mathrm{seen}} = \set{v}$ where $v$ is the current start.
\end{invariant}

\begin{invariant}[Path length and acyclicity]
\label{inv:dfs:path}
Every emitted $\pi \in R$ has $\abs{\pi} = k + 1$, and its nodes are
pairwise distinct.
\end{invariant}

\begin{implnote}
The invariants of Chapter~\ref{chap:dfs} are tested in
\texttt{core/graph-engine/tests/parsers.rs} and
\texttt{core/matcher-ffi/tests/simd\_matcher\_test.rs}. The equivalence
between this DFS matcher and the SIMD implementation of
Chapter~\ref{chap:simd} is pinned by
\texttt{core/matcher-ffi/tests/equivalence\_proptest.rs}.
\end{implnote}

\section{References}
\label{sec:dfs:refs}
\begin{itemize}
  \item T.~H.~Cormen, C.~E.~Leiserson, R.~L.~Rivest, C.~Stein.
        \emph{Introduction to Algorithms}, 3rd ed., \S22.3.
        MIT Press, 2009.~\cite{clrs}
  \item R.~Tarjan. ``Depth-first search and linear graph
        algorithms.'' \emph{SIAM J.\ Comput.}, 1(2), 1972.
  \item W.~Fan, J.~Li, S.~Ma, N.~Tang, Y.~Wu, Y.~Wu. ``Graph Pattern
        Matching: From Intractable to Polynomial Time.''
        \emph{VLDB}, 2010.
\end{itemize}

% ============================================================
\chapter{Path Anomaly Scoring}
\label{chap:scoring}

\section{Problem}
\label{sec:scoring:problem}

\begin{problem}[Weighted five-component anomaly score]
\label{prob:scoring}
Given a finalized graph $G$ augmented with observation counts, a
path $\pi = v_0 \to \dots \to v_k$ produced by
Algorithm~\ref{alg:dfs}, and a non-negative weight vector
$w = (w_1, w_2, w_3, w_4, w_5)$, compute a composite anomaly score
$S(\pi) \in [0, 1]$ and a breakdown
$C(\pi) = (C_1, C_2, C_3, C_4, C_5) \in [0,1]^5$ of the five
contributing components.
\end{problem}

\section{Components}
\label{sec:scoring:components}

Every component is bounded in $[0, 1]$ and averaged over the
appropriate part of the path.

\subsection{Entity rarity}
\label{sec:scoring:er}

For a node $v$ with observation count $\freq(v)$ and
$\freq_{\max} = \max_v \freq(v)$,
\[
\mathrm{ER}(v) \;=\; \max\!\left(0, \; 1 - \frac{\ln \freq(v)}{\ln \freq_{\max}}\right),
\]
with $\mathrm{ER}(v) = 0$ when $\freq_{\max} = 1$ (degenerate case).
Averaged over the path:
\[
C_1(\pi) = \frac{1}{\abs{\pi}} \sum_{v \in \pi} \mathrm{ER}(v).
\]
The logarithmic normalization flattens the heavy-tailed observation
distributions we observe in production so that a single very-popular
node does not crush the score dynamic range.

\subsection{Edge rarity}
\label{sec:scoring:edger}

For a directed pair $(s, d)$ with co-occurrence count
$\freq(s, d)$ and $\freq^{\mathrm{pair}}_{\max}$ the global maximum,
\[
\mathrm{EdgeR}(s, d) \;=\; \max\!\left(0, \; 1 - \frac{\ln \freq(s, d)}{\ln \freq^{\mathrm{pair}}_{\max}}\right),
\]
averaged over the $k = \abs{\pi} - 1$ consecutive pairs in $\pi$.

\subsection{Neighborhood concentration}
\label{sec:scoring:nc}

Precomputed during \textsc{Finalize}$(G)$ and cached per node. The
scalar $\mathrm{NC}(v) \in [0, 1]$ encodes the classical ``dense
anomalous subgraph'' heuristic of Akoglu \emph{et al.}~\cite{akoglu}:
nodes with few neighbors dominated by a single type score high;
nodes embedded in diverse, populous neighborhoods score low.
Averaged over the path.

\subsection{Temporal novelty}
\label{sec:scoring:tn}

Given the first-seen timestamp $\tau_{\mathrm{first}}(v)$ and the
graph-wide extrema $\tau_{\min}, \tau_{\max}$,
\[
\mathrm{TN}(v) \;=\; \frac{\tau_{\mathrm{first}}(v) - \tau_{\min}}{\tau_{\max} - \tau_{\min}},
\]
with $\mathrm{TN}(v) = 0$ when $\tau_{\max} = \tau_{\min}$ (all
nodes seen at the same timestamp). Averaged over the path.

\subsection{GNN threat}
\label{sec:scoring:gnn}

If a graph neural network has produced a threat-probability estimate
$\mathrm{GNN}(v) \in [0, 1]$ per node, the fifth component is the
path average of those probabilities. When the GNN cache is empty
(no inference has run) the component is zero by convention, and the
weight renormalization in \S\ref{sec:scoring:composite} absorbs the
deficit by redistributing mass across the remaining components.

\section{Composite}
\label{sec:scoring:composite}

The composite score is the weight-normalized convex combination
\begin{equation}
\label{eq:composite}
S(\pi) \;=\; \frac{\sum_{i=1}^{5} w_i \cdot C_i(\pi)}{\sum_{i=1}^{5} w_i},
\end{equation}
clamped to $[0, 1]$. When $\sum_i w_i = 0$ the quotient would be
undefined; the implementation returns $S(\pi) = 0$ explicitly in
that branch. This is the invariant that prevents \texttt{NaN} from
propagating into downstream sort comparators, where it would
violate strict-weak ordering.

\section{Invariants}
\label{sec:scoring:invariants}

\begin{invariant}[Boundedness]
\label{inv:scoring:bound}
For every $\pi$ and every non-negative $w$, $S(\pi) \in [0, 1]$ and
$C_i(\pi) \in [0, 1]$ for $i \in \set{1, \dots, 5}$. No component
ever returns $\mathtt{NaN}$ or $\pm\infty$.
\end{invariant}

\begin{invariant}[Zero-weights guard]
\label{inv:scoring:zero}
If $\sum_i w_i = 0$, then $S(\pi) = 0$.
\end{invariant}

\begin{invariant}[Un-finalized scorer]
\label{inv:scoring:unfin}
If \textsc{Finalize}$(G)$ has not been called, \textsc{Score-Path}
returns $(0, \mathbf{0})$ irrespective of $\pi$. Callers use
\textsc{Is-Finalized} as a gate.
\end{invariant}

\begin{theorem}[Entity-rarity monotonicity]
\label{thm:scoring:monotone}
Let $a, b \in V$ satisfy $\freq(a) \le \freq(b)$ and
$\tau_{\mathrm{first}}(a) = \tau_{\mathrm{first}}(b)$. Then
$\mathrm{ER}(a) \ge \mathrm{ER}(b)$.
\end{theorem}

\begin{proof}
$\mathrm{ER}(v) = 1 - \ln \freq(v) / \ln \freq_{\max}$ is
non-increasing in $\freq(v)$ because $\ln$ is non-decreasing and
$\ln \freq_{\max} > 0$ whenever $\freq_{\max} > 1$. The
$\max(0, \cdot)$ clamp preserves the ordering.
\end{proof}

\begin{implnote}
Invariants~\ref{inv:scoring:bound}--\ref{inv:scoring:unfin} and
Theorem~\ref{thm:scoring:monotone} are pinned by the property tests
in \texttt{core/graph-engine/tests/anomaly\_proptest.rs}, running
256 cases per property across randomized weight vectors and
observation distributions.
\end{implnote}

\section{Trade-offs}
\label{sec:scoring:tradeoffs}

\begin{itemize}
  \item \textbf{Averaging masks outliers.} A single extremely-rare
        node on a long path gets averaged down by its common
        neighbors. Alternatives (max-over-path, quantile) break the
        monotonicity property and complicate paginated consumers.
  \item \textbf{Log normalization assumes Zipfian distributions.}
        If observation counts are roughly uniform, $C_1$ and $C_2$
        collapse to near zero.
  \item \textbf{First-seen is spoofable.} An attacker aware of the
        scoring can backdate $\tau_{\mathrm{first}}(v)$ on a fresh
        node to suppress $C_4$. Mitigation is out of scope; ingestion
        must cross-check against an append-only audit log.
\end{itemize}

\section{References}
\label{sec:scoring:refs}
\begin{itemize}
  \item L.~Akoglu, H.~Tong, D.~Koutra. ``Graph-based anomaly detection
        and description: a survey.'' \emph{DMKD} 29(3), 2015.~\cite{akoglu}
  \item V.~Chandola, A.~Banerjee, V.~Kumar. ``Anomaly detection: a
        survey.'' \emph{ACM Computing Surveys} 41(3), 2009.
  \item L.~Akoglu, M.~McGlohon, C.~Faloutsos. ``OddBall: Spotting
        Anomalies in Weighted Graphs.'' \emph{PAKDD}, 2010.
\end{itemize}

% ============================================================
\chapter{SIMD-Accelerated Sorted Set Intersection}
\label{chap:simd}

\section{Problem}
\label{sec:simd:problem}

\begin{problem}[Sorted set intersection]
\label{prob:simd}
Given two strictly ascending arrays $A[0 \mathinner{\ldotp\ldotp} m)$ and
$B[0 \mathinner{\ldotp\ldotp} n)$ of 32-bit unsigned integers, compute either
\begin{itemize}[leftmargin=2em]
  \item $\textsc{Intersect}(A, B, \mathit{out})$: write the sorted
        intersection $A \cap B$ into $\mathit{out}$ and return
        $\abs{A \cap B}$; or
  \item $\textsc{Intersect-Count}(A, B)$: return $\abs{A \cap B}$
        without materializing the output (used when the caller
        only needs the count for pruning).
\end{itemize}
\end{problem}

A temporal variant \textsc{Intersect-Temporal} additionally filters
matches whose associated timestamp lies outside a window
$[t_{\text{lo}}, t_{\text{hi}}]$. Used by the matcher of
Chapter~\ref{chap:dfs} when narrowing candidate edges.

\section{Dispatch}
\label{sec:simd:dispatch}

A static feature-detection routine selects one of four variants at
first call and memoizes the choice in an atomic global:

\begin{table}[h]
\centering
\small
\begin{tabularx}{\linewidth}{l l l X}
\toprule
\textbf{ISA} & \textbf{Variant} & \textbf{Lane width} & \textbf{Intrinsic family} \\
\midrule
AVX-512 & scalar fallback (for now) & --- & --- \\
AVX2    & \textsc{Intersect-Avx2}    & $8 \times \code{uint32}$ & \code{\_mm256\_cmpeq\_epi32}, \code{\_mm256\_movemask\_epi8} \\
SSE4.2  & scalar fallback            & --- & --- \\
NEON    & scalar fallback            & --- & --- \\
Other   & \textsc{Intersect-Scalar}  & 1 & --- \\
\bottomrule
\end{tabularx}
\end{table}

Branch prediction eliminates the dispatch cost after the first few
invocations.

\section{Scalar variant}
\label{sec:simd:scalar}

The two-pointer merge, for reference.

\begin{algorithm}[H]
\caption{\textsc{Intersect-Scalar}$(A, B, \mathit{out})$}
\label{alg:simd:scalar}
\Input{sorted arrays $A[0 \mathinner{\ldotp\ldotp} m)$, $B[0 \mathinner{\ldotp\ldotp} n)$}
\Output{count $k = \abs{A \cap B}$; $\mathit{out}[0 \mathinner{\ldotp\ldotp} k)$ holds the intersection}
$i \gets 0$; $j \gets 0$; $k \gets 0$\;
\While{$i < m$ \KwAnd $j < n$}{
  \lIf{$A[i] = B[j]$}{$\mathit{out}[k] \gets A[i]$; $k \gets k+1$; $i \gets i+1$; $j \gets j+1$}
  \lElseIf{$A[i] < B[j]$}{$i \gets i+1$}
  \lElse{$j \gets j+1$}
}
\KwReturn $k$\;
\end{algorithm}

Running time is $\Theta(m + n)$ in both the best and worst cases.

\section{AVX2 variant}
\label{sec:simd:avx2}

The AVX2 specialization follows the broadcast--compare--movemask
scheme of Schlegel, Gubichev, and Neumann~\cite{schlegel}. Let
$V = 8$ be the AVX2 lane count.

\begin{algorithm}[H]
\caption{\textsc{Intersect-Avx2}$(A, B, \mathit{out})$}
\label{alg:simd:avx2}
\Input{sorted arrays $A[0 \mathinner{\ldotp\ldotp} m)$, $B[0 \mathinner{\ldotp\ldotp} n)$, with strictly increasing values}
\Output{count $k$; $\mathit{out}[0 \mathinner{\ldotp\ldotp} k)$ holds the intersection}
$i \gets 0$; $j \gets 0$; $k \gets 0$\;
\While{$i + V \le m$ \KwAnd $j < n$}{
  $\mathbf{v}_A \gets \textsc{Load-256}(A + i)$ \Comment*[r]{8 lanes of $A$}
  $a_{\max} \gets A[i + V - 1]$\;
  \While{$j < n$ \KwAnd $B[j] \le a_{\max}$}{
    $\mathbf{v}_B \gets \textsc{Broadcast-256}(B[j])$\;
    $\mathbf{c} \gets \mathbf{v}_A \stackrel{?}{=} \mathbf{v}_B$ \Comment*[r]{lane-wise equality}
    $M \gets \textsc{Movemask}(\mathbf{c})$\;
    \lIf{$M \ne 0$}{$\mathit{out}[k] \gets B[j]$; $k \gets k + 1$}
    $j \gets j + 1$\;
  }
  $i \gets i + V$\;
}
\textbf{scalar tail} for $i \ge m - V + 1$ \;
\KwReturn $k$\;
\end{algorithm}

Because both arrays are strictly ascending, at most one lane of
$\mathbf{c}$ can be nonzero per broadcast; the $\textsc{Movemask}$
result therefore encodes a simple ``hit or miss.'' The scalar tail
ensures no out-of-bounds load when $m \not\equiv 0 \pmod{V}$.

\section{Complexity}
\label{sec:simd:complexity}

Let $m, n$ be the input sizes and $k = \abs{A \cap B}$ the output
size.

\begin{theorem}[AVX2 asymptotics]
\label{thm:simd:avx2}
\textsc{Intersect-Avx2} runs in $\bigO(m + n)$ time. Empirically,
the constant factor is $3$--$4\times$ smaller than
\textsc{Intersect-Scalar} on uniform-density inputs, and up to
$6$--$8\times$ smaller when one array is sparse relative to the
other (the inner $B$-loop executes multiple broadcasts per outer
$A$-step).
\end{theorem}

\begin{remark}[Pathological case]
When every $B[j] < \max(\mathbf{v}_A)$ for nearly every outer step,
the inner loop can traverse the full length of $B$ per outer
iteration, degenerating toward $\bigO(m \cdot n)$. The matcher
sidesteps this by invoking intersection with the smaller array as
the outer loop.
\end{remark}

\section{Invariants}
\label{sec:simd:invariants}

\begin{invariant}[Input well-formedness]
\label{inv:simd:sorted}
$A$ and $B$ are each strictly ascending with no duplicates.
Violation causes silent duplicates in the output; the matcher's
adjacency-list builder guarantees this property upstream.
\end{invariant}

\begin{invariant}[Output capacity]
\label{inv:simd:capacity}
The caller allocates $\mathit{out}$ with room for at least
$\min(m, n)$ elements.
\end{invariant}

\begin{invariant}[No SIMD overread]
\label{inv:simd:bounds}
The vectorized loop processes only full $V$-lane vectors; the tail
falls back to scalar. No unaligned \code{loadu} reads past the end
of $A$ or $B$, so the algorithm never takes a page fault beyond the
allocated regions.
\end{invariant}

\begin{invariant}[Thread-safety]
\label{inv:simd:thread}
The ISA selection is written once behind an atomic flag; the
specialized variants are pure functions. Safe to call concurrently
from multiple matcher threads.
\end{invariant}

\begin{invariant}[ABI stability]
\label{inv:simd:abi}
The Rust bridge asserts \code{gm\_abi\_version() == EXPECTED\_ABI\_VERSION}
at first call via \code{std::sync::Once}. Silent drift in the
dispatch table layout therefore fails loudly at initialization
rather than at runtime with memory corruption.
\end{invariant}

\begin{implnote}
The full SIMD matcher is tested end-to-end (not intersection alone)
by \texttt{core/matcher-ffi/tests/equivalence\_proptest.rs}, which
enumerates random small graphs and hypotheses and requires the DFS
and SIMD result sets to be identical. A bug in
\textsc{Intersect-Avx2} surfaces there as a missing path.
\end{implnote}

\section{References}
\label{sec:simd:refs}
\begin{itemize}
  \item B.~Schlegel, A.~Gubichev, T.~Neumann. ``Fast Sorted-Set
        Intersection using SIMD Instructions.'' \emph{ADMS},
        2011.~\cite{schlegel}
  \item D.~Lemire, O.~Kurz, C.~Rupp. ``Fast random integer generation
        in an interval.'' \emph{ACM TOMS} 45(1), 2019.
  \item \emph{Intel 64 and IA-32 Architectures Software Developer's
        Manual}, Vol.\ 2B, \S\code{VPCMPEQD}, \code{VPMOVMSKB}.
\end{itemize}

% ============================================================
\backmatter

\chapter*{Bibliography}
\addcontentsline{toc}{chapter}{Bibliography}
\begin{thebibliography}{9}
  \bibitem{clrs}
    T.~H.~Cormen, C.~E.~Leiserson, R.~L.~Rivest, C.~Stein.
    \emph{Introduction to Algorithms}, 3rd ed. MIT Press, 2009.

  \bibitem{akoglu}
    L.~Akoglu, H.~Tong, D.~Koutra.
    ``Graph-based anomaly detection and description: a survey.''
    \emph{Data Mining and Knowledge Discovery}, 29(3):626--688, 2015.

  \bibitem{schlegel}
    B.~Schlegel, A.~Gubichev, T.~Neumann.
    ``Fast Sorted-Set Intersection using SIMD Instructions.''
    \emph{International Workshop on Accelerating Data Management
    Systems Using Modern Processor and Storage Architectures
    (ADMS)}, 2011.
\end{thebibliography}

\end{document}
